\documentclass[12pt]{article}
\usepackage[utf8]{inputenc}
\usepackage[T1]{fontenc}
\usepackage{hyperref}
\title{Thompson Language Resources for K-12 Curriculum Development}
\author{Compiled by Claude Sonnet 5 (claude-sonnet-5) for Andr\'as Kornai}
\date{}
\begin{document}
\maketitle

\section{Introduction}
Thompson (Nlaka'pamux, \textit{nłeʔkepmxcín}; ISO 639-3 thp) is a Salishan language of British Columbia, Canada. Sources report about 130 first-language speakers (Ethnologue 2016, as cited by the Endangered Languages Project), and other sources give higher figures. This report is based on a small number of web searches. I found no ready-made K-12 curriculum.

\section{Primary Educational Resources}
I found no textbooks, primers or graded readers for Thompson in the search results. The Nlaka'pamux Tribal Council in Lytton is reported to run a language teaching programme, but I saw no curriculum document. The FirstVoices dictionary below has teaching features and could support classroom work.

\section{Linguistic and Reference Materials}
\begin{itemize}
\item Thompson River Salish Dictionary (Laurence and Terry Thompson, 1996, UMOPL 12). It is a bilingual dictionary in both directions and covers eight dialect areas. \url{https://wiki.ubc.ca/Documentation:RelLex/Thompson_River_Salish_Dictionary}
\item UBC Arts Knowledge Base entry for the same dictionary. \url{https://knowledgebase.arts.ubc.ca/?p=3264}
\item Omniglot overview of the language and its writing systems. \url{https://www.omniglot.com/writing/thompson.htm}
\end{itemize}

\section{Religious and Community Resources}
No scripture translations were found in the search results. A listing of Interior Salish language resources is at \url{https://TOTABC.org/ind/resource/languages}.

\section{Digital and Audio Resources}
\begin{itemize}
\item The nłeʔkepmxcin FirstVoices dictionary is maintained by the Scw'exmx Tribal Council. The record says it was last updated in August 2023. \url{https://wiki.ubc.ca/Documentation:RelLex/n%C5%82e%CA%94kepmxcin_FirstVoices}
\item Endangered Languages Project entry. \url{https://endangeredlanguages.com/node/85273}
\end{itemize}

\section{Curriculum Development Context}
Few resources exist, and most are lexical. A curriculum would need new material built from the dictionary and FirstVoices audio, developed with the Nlaka'pamux community and its language programmes. The community orthography in use should be confirmed before materials are made, because Americanist and Bouchard orthographies are both in use.
\end{document}